% \newpage
\section{Discussion}
\label{sec:summary}

We introduced an open-source library, \libName, that incorporates the paradigms of conversable agents and conversation programming. This library utilizes capable agents that are well-suited for multi-agent cooperation. It features a unified conversation interface among the agents, along with an auto-reply mechanisms, which help establish an agent-interaction interface that capitalizes on the strengths of chat-optimized LLMs with broad capabilities while accommodating a wide range of applications. \libName serves as a general framework for creating and experimenting with multi-agent systems that can easily fulfill various practical requirements, such as reusing, customizing, and extending existing agents, as well as programming conversations between them.

Our experiments, as detailed in Section~\ref{sec:app}, demonstrate that this approach offers numerous benefits. 
The adoption of \libName has resulted in improved performance (over state-of-the-art approaches), reduced development code, and decreased manual burden for existing applications. It offers flexibility to developers, as demonstrated in A1 (scenario 3), A5, and A6, where \libName enables multi-agent chats to follow a dynamic pattern rather than fixed back-and-forth interactions. It allows humans to engage in activities alongside multiple AI agents in a conversational manner. Despite the complexity of these applications (most involving more than two agents or dynamic multi-turn agent cooperation), the implementation based on \libName remains straightforward. Dividing tasks among separate agents promotes modularity. 
Furthermore, since each agent can be developed, tested, and maintained separately, this approach simplifies overall development and code management.

Although this work is still in its early experimental stages, it paves the way for numerous future directions and research opportunities. For instance, we can explore effective integration of existing agent implementations into our multi-agent framework and investigate the optimal balance between automation and human control in multi-agent workflows. As we further develop and refine \libName, we aim to investigate which strategies, such as agent topology and conversation patterns, lead to the most effective multi-agent conversations while optimizing the overall efficiency, among other factors. While increasing the number of agents and other degrees of freedom presents opportunities for tackling more complex problems, it may also introduce new safety challenges that require additional studies and careful consideration.

We provide more discussion in Appendix~\ref{append:summary}, including guidelines for using \libName and direction of future work.
We hope \libName will help improve many LLM applications in terms of speed of development, ease of experimentation, and overall effectiveness and safety. We actively welcome contributions from the broader community. 

